\documentclass[10pt,a4paper]{amsart}
\usepackage[margin=19mm]{geometry}
\usepackage{amsmath,amssymb,mathtools}
\usepackage[scr=rsfs,cal=euler]{mathalfa}
\usepackage{xfrac}
\usepackage[unicode,hidelinks,bookmarks=false]{hyperref}
\newcommand{\ie}{i.e.\ }
\newcommand{\cf}{cf.\ }
\newcommand{\resp}{respectively\ }
\newcommand{\Subsection}{\subsection}
\providecommand{\nobreakdash}{\nobreak\hbox{-}}
\setlength{\emergencystretch}{2em}
\multlinegap0pt
\usepackage[shortlabels]{enumitem}
\usepackage{amssymb}
\usepackage{dsfont}
\usepackage{float}
\usepackage{bm}
\usepackage{algorithm}\usepackage{algpseudocode}
\usepackage[normalem]{ulem}
\usepackage{xcolor}
\newtheorem{prop}{Proposition}
\newcommand{\hj}{h_{j\vect{\beta}}}
\newcommand{\phj}{\phi_{j\vect{\gamma}}}
\newcommand{\thj}{\widetilde{h}_{j\vect{\beta}}}
\newcommand{\tphj}{\widetilde{\phi}_{j\vect{\gamma}}}
\newcommand{\vect}[1]{\boldsymbol{#1}}
\title{\bf A New Class of Asymptotically Distribution-Free Smooth Tests}
\author{Xiangyu Zhang, Sara Algeri}
\date{Source: arXiv:2508.01973v4 (CC BY 4.0)}
\begin{document}
\maketitle

\noindent\emph{Source and licence.} \bf A New Class of Asymptotically Distribution-Free Smooth Tests. Version arXiv:2508.01973v4 (CC BY 4.0), distributed by arXiv under the Creative Commons Attribution 4.0 International licence, http://creativecommons.org/licenses/by/4.0/, as recorded in the arXiv licence metadata for this exact version and shown on the abstract page https://arxiv.org/abs/2508.01973v4. Full licence notice: \texttt{licenses/arxiv-2508.01973-CC-BY-4.0.txt}.\par\smallskip

\noindent\textit{Editorial scope.} Counted unit: the article's prop prop:6 (source0.tex lines 638--645). The article's complete original proof of it is delivered with it (source0.tex lines 646--649, one \texttt{proof} environment of 4 lines). No mathematics is written, repaired or re-derived: the statement and the proof are the article's own lines, in the article's own notation, and no retained line is substituted or re-flowed.  No further source-local context is needed: every cross-reference of every retained line resolves inside the two delivered spans.  Nothing was omitted from any retained span, so the package carries no omission record.  No retained line contains a citation, so the wrapper carries no bibliography and no bibliography entry.  No retained line contains a figure environment, an image-inclusion command or drawing code, so no image is delivered and the package declares no asset dependency.  Editorial formatting only, in addition: an AMS \texttt{amsart} preamble that restates the article's own declarations and macros used by the retained lines, namely the environments \texttt{prop} on separate counters, together with the standard packages those lines need.  The retained environments print in this document as Proposition 1 in order of appearance, whereas the article prints the same items with the numbering of its own full document.  The complete native source of the article as distributed in the arXiv e-print for this exact version is bundled unchanged as \texttt{sources/182648/source0.tex}.
\begin{prop} \label{prop:6}
The set of functions $\{\vect{U}_{p}K l_{\vect{\gamma},\vect{\beta}} \phj\}_{j=1}^\infty$ forms an orthonormal basis for $L^2(G_{\vect{\beta}})$. 
Define, for each $j \geq 1, $ 
\begin{equation*}
\hj = \vect{U}_{p}K l_{\vect{\gamma},\vect{\beta}} \phj, \quad \text{and} \quad \thj = \vect{U}_{p}K l_{\vect{\gamma},\vect{\beta}} \tphj.
\end{equation*}
Then, the empirical processes $v_{\vect{\beta},n}(\thj)$ and $v_{\vect{\gamma},n}(\tphj)$ have the same limiting distribution under $G_{\vect{\beta}}$ and $F_{\vect{\gamma}}$, respectively.
\end{prop} 

\begin{proof}
The first claim follows from the unitary nature of the operators $\vect{U}_{p}$, $K$, and $l_{\vect{\gamma},\vect{\beta}}$.  
The second claim is a direct consequence of the asymptotic Gaussianity of $v_{\vect{\beta},n}(\thj)$ and $v_{\vect{\gamma},n}(\tphj)$ and the construction of the $\{\thj\}_{j=1}^\infty$ basis, which ensures that the empirical process indexed by such functions has, under $G_{\vect{\beta}}$, the same mean and covariance as the empirical process indexed by functions $\{\tphj\}_{j=1}^\infty$ under $F_{\vect{\gamma}}$.
\end{proof}

\end{document}
